\section{An exact representation and comparison at boundary incidences}
\label{sec:comparison}

The bulk comparison of Section~\ref{sec:stability} concerns unconditioned
torus laws.  We now obtain the two tools needed to use it in a domain with
prescribed exterior spins.  The first orders spins across certain
\emph{strict} changes of boundary incidences.  The second compares all
interior spin events, pointwise in their densities, across a fixed free
buffer.  Neither tool requires positive association of the perturbed spin
law with arbitrary pins.

Throughout this section the inverse temperature is the value supplied by
Theorem~\ref{thm:critical-trajectory}.  We abbreviate the conditional spin
law of Lemma~\ref{lem:contour-spin} by
\(\nu_{\Omega,\lambda}^{\tau}
=\nu_{\Omega,\beta_c(J,U,\lambda),\lambda}^{\tau}\),
and write \(\mu_{\lambda,N}\) for the unconditioned law on a square torus
of side \(N\).  We take tori larger than the fixed interaction range and
than all the charts used below.  All thresholds in the local construction
depend only on the fixed interaction, not on a domain or its pins.

\subsection{Private auxiliary variables for a local factor}

Write \(K_0=\beta_0J\).  The spin weight can be expressed as
\begin{equation}
 \exp\left(K_0\sum_{uv}\sigma_u\sigma_v\right)
 \prod_{\alpha,x}F_\alpha(\sigma_{x+A_\alpha}),
 \qquad
 \max_\alpha\|F_\alpha-1\|_\infty\longrightarrow0
 \quad(\lambda\longrightarrow0).
 \label{eq:comparison-local-factors}
\end{equation}
Here the sum is over nearest-neighbor bonds, the index \(\alpha\) ranges
over a fixed finite list of finite supports, and \(x\) ranges over their
translates.  Every \(F_\alpha\) is positive.  Lemma~\ref{lem:contour-spin}
gives the factors coming from \(U\); the change
\(\beta_cJ-K_0\) is included as another bond factor.  Enlarging a support
does not change its factor.  Thus we may fix, for each \(\alpha\), a finite
nearest-neighbor tree \(T_\alpha\) whose vertex set contains \(A_\alpha\).
Singleton supports are allowed.

We use binary variables with values in \(\{-1,+1\}\).  A variable placed
on a bond will be called a \emph{bit}, to distinguish it from a physical
spin at a lattice vertex.  For a fixed integer \(M\), put \(M\) bits
\(b_{e,1},\ldots,b_{e,M}\) on every nearest-neighbor bond \(e=uv\).
Choose \(a>0\) so that
\begin{equation}
 \frac{M}{2}\log\cosh(2a)=K_0.
 \label{eq:comparison-bit-coupling}
\end{equation}
Then summing the bits in
\(\prod_{i=1}^M\exp(a(\sigma_u+\sigma_v)b_{e,i})\)
gives a spin-independent constant times
\(\exp(K_0\sigma_u\sigma_v)\).  In particular, every bit has a strictly
positive coupling to both ends of its bond.

The reason for using these variables is a strict distinction between
free spins and pins.  At the unperturbed vertex, summing the physical
spin gives $2\cosh(a\sum_i b_i)$, whereas fixing it to either sign
gives $\exp(\pm a\sum_i b_i)$.  In the first factor the conditional
odds of one bit increase with every other incident bit.  Its logarithmic
increments lie strictly between those of the minus and plus factors.
We will retain these finite-table inequalities under perturbation.
The construction below must therefore realize each local factor exactly
while changing every vertex table by a quantity tending to zero.

\begin{proposition}[Exact vertex representation]
\label{prop:bit-representation}
There is a finite choice of \(M\) such that, for all sufficiently small
\(\lvert\lambda\rvert\), positive functions
\[
 W_v\bigl(\sigma_v,b_{\mathcal B(v)}\bigr)
\]
represent the spin weight in \eqref{eq:comparison-local-factors} exactly,
up to a constant, as the spin marginal of \(\prod_v W_v\).
Here \(\mathcal B(v)\) is the set of the \(4M\) bits on bonds incident to
\(v\).  The same functions are used whether a physical spin is free or
pinned.  There are only finitely many local tables, and uniformly in
their entries
\begin{equation}
 W_v(\sigma,b)\longrightarrow
 W_v^0(\sigma,b):=\exp\left(a\sigma\sum_{i\in\mathcal B(v)}b_i\right)
 \quad\text{as }\lambda\longrightarrow0.
 \label{eq:comparison-table-limit}
\end{equation}
The representation applies to either sign of every sufficiently small
local perturbation.
\end{proposition}

\begin{proof}
We construct one factor on a tree \(T\), and then give different
constructions disjoint sets of bits.  For each complete spin pattern
\(t\in\{-1,+1\}^{V(T)}\), reserve two private bits on every edge of
\(T\), with the same coupling \(a\).  Conditional on all the physical
spins, these bits are independent in the unperturbed representation.
For the two bits on \(e=uv\), set \(D_e=b_{e,1}-b_{e,2}\).  Their
conditional moments are
\begin{equation}
 \E[D_e\mid\sigma]=0,\qquad
 q_e(\sigma_u,\sigma_v)
 :=\E[D_e^2\mid\sigma]
 =2\operatorname{sech}^2\bigl(a(\sigma_u+\sigma_v)\bigr)>0.
 \label{eq:comparison-duplicate-moments}
\end{equation}
At a vertex \(v\) of this private copy of \(T\), multiply its vertex
factor by
\begin{equation}
 1+c_{v,t}\,\one_{\{\sigma_v=t_v\}}
       \prod_{\substack{e\in E(T)\\e\ni v}}D_e.
 \label{eq:comparison-tree-correction}
\end{equation}
An empty product is one.  In the expansion of the product over vertices,
a chosen subset \(S\subseteq V(T)\) contributes one copy of \(D_e\) for
each end of \(e\) in \(S\).  If \(S\) is nonempty and proper, the
connectedness of \(T\) gives an edge with exactly one end in \(S\).
Its centered difference makes this term vanish upon summing the private
bits.  The only surviving terms are therefore the empty selection and
the full selection.  Their sum is
\begin{equation}
 1+\one_{\{\sigma_{V(T)}=t\}}
       \left(\prod_{v\in V(T)}c_{v,t}\right)
       \left(\prod_{uv\in E(T)}q_{uv}(t_u,t_v)\right).
 \label{eq:comparison-tree-identity}
\end{equation}
The same formula holds for a singleton tree, with an empty edge product.

For a target factor \(F\) on a subset of \(V(T)\), put
\[
 \varepsilon_t=F(t)-1,\qquad
 Q_t=\prod_{uv\in E(T)}q_{uv}(t_u,t_v)>0.
\]
Choose coefficients whose product is \(\varepsilon_t/Q_t\).
For example, if \(m=|V(T)|\), take absolute values
\(\lvert\varepsilon_t/Q_t\rvert^{1/m}\) at every vertex, place the sign
of \(\varepsilon_t\) at one prescribed vertex, and take all coefficients
zero when \(\varepsilon_t=0\).  Every coefficient tends to zero as
\(F\to1\), for either sign of \(\varepsilon_t\).
Equation~\eqref{eq:comparison-tree-identity} then realizes
\(1+\varepsilon_t\one_{\{\sigma_{V(T)}=t\}}\).
Use distinct private bits for different patterns.  Conditional
independence allows their contributions to be multiplied, and exactly
one pattern is active at a given physical spin configuration.  Hence
\[
 \prod_t\left(1+\varepsilon_t
                    \one_{\{\sigma_{V(T)}=t\}}\right)=F(\sigma).
\]

There are finitely many trees, finitely many patterns on each, and only
finitely many translated trees using a given bond.  First count these
uses, and choose \(M\) large enough to assign all their pairs of private
bits disjointly.  Unused bits retain their original factor.  This
allocation can be made translation covariant; its number is independent
of the size of the torus.  Choose \(a\) only after this finite allocation,
using \eqref{eq:comparison-bit-coupling}.
The product of all the factors
\eqref{eq:comparison-tree-correction} at \(v\), times \(W_v^0\), defines
\(W_v\).  Each \(D_e\) has absolute value at most two.  Since a vertex
belongs to only finitely many assigned trees, all correction factors are
positive when their coefficients are small enough.  Their tables then
converge uniformly to one, proving
\eqref{eq:comparison-table-limit}.  Summing all private channels, with
the physical spins fixed, proves the claimed exact identity.
\end{proof}

The construction need not be analytic in \(\lambda\): taking a root of
\(\lvert\varepsilon_t\rvert\) provides the required continuity.  The
analytic construction of the critical temperature was completed in
Section~\ref{sec:flow} before this representation was introduced.

\subsection{Strict local margins and their comparison consequence}

For a positive function \(H\) of bits, let
\[
 d_i\log H(b_{-i})
 =\log H(b_i=+1,b_{-i})-\log H(b_i=-1,b_{-i}).
\]
The function is \emph{log-supermodular} if every such increment is
nondecreasing in every other coordinate.  On a Boolean cube this is
equivalent to
\(H(b\wedge b')H(b\vee b')\ge H(b)H(b')\); it suffices to sum the
two-coordinate differences along coordinate paths.  A positive
probability mass with this property is also called MTP\(_2\).

Summing or fixing the spin at \(v\) gives respectively the factors
\[
 F_v(b)=W_v(+1,b)+W_v(-1,b),\qquad
 P_v^+(b)=W_v(+1,b),\qquad P_v^-(b)=W_v(-1,b).
\]
The conditional probability of a free physical spin being plus is
\[
 k_v(b)=\frac{W_v(+1,b)}{F_v(b)}.
\]
We distinguish these vertex factors from the conditional odds of a bit
in the entire measure, which also involve its other endpoint.

\begin{lemma}[Finite-table margins]
\label{lem:comparison-margins}
After decreasing the allowed interval of \(\lambda\), every free factor
\(F_v\) is strictly log-supermodular in each pair of distinct incident
bits, and every kernel \(k_v\) is strictly increasing in each incident
bit.  For every incident coordinate \(i\), moreover,
\begin{equation}
 \sup_b d_i\log P_v^-(b)
 <\inf_b d_i\log F_v(b)
 \le\sup_b d_i\log F_v(b)
 <\inf_b d_i\log P_v^+(b).
 \label{eq:comparison-strict-margins}
\end{equation}
The gaps are bounded below by a positive constant independent of the
volume and the pins.
\end{lemma}

\begin{proof}
At \(\lambda=0\), put \(H=a\sum_{j\ne i}b_j\).  Then
\[
 F_v^0(b)=2\cosh\left(a\sum_jb_j\right),\qquad
 d_i\log F_v^0=\log\cosh(H+a)-\log\cosh(H-a).
\]
Because \(-1<\tanh u<1\) for finite \(u\),
\begin{equation}
 -2a<d_i\log F_v^0<2a.
 \label{eq:comparison-base-pin-gaps}
\end{equation}
The increments of \(\log P_v^{-,0}\) and \(\log P_v^{+,0}\) are
exactly \(-2a\) and \(2a\).  For two distinct bits the mixed difference
of \(\log F_v^0\), writing \(H\) now for the contribution from all the
remaining bits, is
\[
 \log\cosh(H+2a)+\log\cosh(H-2a)-2\log\cosh H>0,
\]
by strict convexity.  Finally, the log odds of the free spin are
\(2a\sum_jb_j\), with increment \(4a>0\) in each bit.
There are only finitely many entries in these tables.  Their strictly
positive minimum gaps persist under the uniform positive perturbation
in \eqref{eq:comparison-table-limit}.  This proves all the assertions.
\end{proof}

The free factors therefore give an attractive law on bits when all
physical spins are free.  The factors of a pinned spin have not been
claimed attractive.  The strict inequalities
\eqref{eq:comparison-strict-margins} instead provide exactly the
cross-law comparisons needed below.  The comparison uses the
order-preserving Markov-chain argument underlying Holley's theorem
\citep[Theorem~(6)]{Holley1974}; here the strict incidence conditions
will order the single-bit conditional probabilities.

\begin{proposition}[Comparison at strict incidences]
\label{prop:strict-incidence}
Consider two laws formed from the same vertex factors \(W_v\), with
possibly different sets of pinned physical spins.  Label the laws lower
and upper.  Let \(V\) be a union of nearest-neighbor connected components
of their common free vertex set.  At each vertex outside \(V\) adjacent
to \(V\), require one of the following possibilities:
\begin{center}
\begin{tabular}{c|c}
lower law & upper law\\ \hline
free & pinned plus\\
pinned minus & free\\
pinned minus & pinned plus
\end{tabular}
\end{center}
Then there is a coupling in which all bits on bonds touching \(V\), and
all physical spins in \(V\), are lower than or equal to their upper-law
counterparts.  The assertion remains valid if arbitrary values of bits
on bonds not touching \(V\) have been conditioned on separately in the
two laws.  It consequently remains valid for arbitrary mixtures of
these exterior bit conditions.
\end{proposition}

\begin{proof}
Let \(\mathcal B(V)\) consist of all bits on bonds with at least one
endpoint in \(V\).  Condition in each law on all other bits.  After
summing every free physical spin, the conditional weight on
\(\mathcal B(V)\) is a product of vertex factors, restricted to the
conditioned values where necessary.  A vertex in \(V\) has every
incident bit in \(\mathcal B(V)\).  Its factor is the identical
attractive function \(F_v\) in the two laws.  Its contribution to the
conditional log odds of a bit is therefore ordered whenever the two
configurations of varying bits are ordered.

At an endpoint outside \(V\), the table in the statement and
\eqref{eq:comparison-strict-margins} order the two contributions to those
log odds, regardless of the values of all other incident bits.  This
last qualification allows both the separately conditioned bits and the
other varying bits at that endpoint to differ.  Adding the contributions
of the two endpoints orders the full single-bit conditional log odds.

Run single-bit heat-bath chains in the two finite state spaces, updating
the same coordinate and using the same uniform random number at each
step.  The ordered conditional probabilities preserve the coordinatewise
order.  Start from any ordered pair, and take a subsequential stationary
limit of the coupled chains.  Strict positivity makes each marginal
chain irreducible and gives its prescribed Gibbs measure as its unique
stationary law.  We have thus obtained an ordered coupling of the bits.

Conditional on all bits, the free physical spins are independent, with
the kernels \(k_v\).  At the vertices of \(V\) these kernels are the
same in both laws and are increasing by
Lemma~\ref{lem:comparison-margins}.  Another family of common uniform
random numbers gives the ordered physical spins.  Finally, the argument
worked for every pair of exterior bit configurations.  Integrating these
couplings against the two exterior distributions proves the mixture
assertion.
\end{proof}

Two equal-sign pins at a boundary vertex are deliberately absent from
the table.  Their perturbed factors need not have ordered increments
when their other incident bits differ.  In particular,
Proposition~\ref{prop:strict-incidence} does not imply unrestricted
boundary monotonicity or the FKG inequality for pinned physical spins.

\subsection{Stopped reveals and geometric shielding}

The spin paths used here lie in the fixed triangulation obtained by
adding the northwest--southeast diagonal to every elementary square.
On bulk faces this realizes the contour pairing in
Section~\ref{sec:introduction}.  Its edges do not cross, and it contains
every nearest-neighbor bond.  We record the precise role of a reveal, so
that conditioning on a crossing does not introduce an unexamined
assumption about a spatial Markov property of the spin interaction.

\begin{lemma}[Reveals with a monochromatic facing boundary]
\label{lem:comparison-stopped-reveal}
Suppose spins are queried one at a time, each next query and the stopping
decision being determined by the preceding queries and their answers.
Conditional on a complete transcript, the remaining law is exactly the
original law with the revealed spins pinned.  This remains true after
conditioning on any event determined by the transcript.

In a finite triangulated disk with a prescribed inner disk, one can
explore from its outer boundary for a shield of either color \(s\).
On success, no inner-disk spin has been queried, and every queried
neighbor of the remaining component containing that disk has color
\(s\).  A color-\(s\) circuit in a narrower annular band separating the
two disks guarantees success.  All queried sites are on the exterior
side of that circuit or on the circuit itself.

The corresponding exploration in a triangulated topological rectangle,
from one side toward the opposite side, determines whether there is a
color-\(s\) crossing between the other two sides.  On success its facing
boundary contains such a crossing, and the exploration queries no spins
strictly beyond it toward the opposite side.  If that crossing extends
past the ends of a comparison region, every nearest-neighbor incidence
from the region across the explored side meets a pin of color \(s\).
\end{lemma}

\begin{proof}
Fix a deterministic query order to resolve all choices.  A transcript
lists queried vertices \(v_1,\ldots,v_m\) and their values
\(s_1,\ldots,s_m\).  Any configuration agreeing with these values causes
the same successive choices and the same stopping decision, by induction
on the number of queries.  Thus the transcript event is the cylinder
\(\{\sigma_{v_i}=s_i:1\le i\le m\}\).  Conditioning on a
transcript-measurable event adds nothing once a compatible transcript
has been specified.  Independent exploration randomness, if used, can
first be conditioned on and then averaged out.

For the annular construction, query the outer graph boundary.  Keep a
queue of all revealed vertices of color \(-s\), and query every
previously unqueried triangulation-neighbor of a queued vertex before
removing that vertex from the queue.  Newly discovered vertices of color
\(-s\) join the queue; vertices of color \(s\) do not.  Declare failure
just before a query would enter the protected inner disk.  Declare
success when the queue has been exhausted.  This is a finite adaptive
procedure of the kind just considered.

On success let \(V\) be the unqueried nearest-neighbor component
containing the inner disk, or the union of such components if the chosen
inner set has more than one component.  An unqueried nearest neighbor
of \(V\) belongs to \(V\).  A queried neighbor cannot have color \(-s\),
since every neighbor of such a vertex was queried before its queue entry
was removed.  Hence every boundary incidence of \(V\) meets a revealed
\(s\) vertex.  If a color-\(s\) circuit separates the protected disk from
the outer boundary, every queued vertex has a color-\(-s\) path to the
outer boundary and so lies outside the circuit.  A further queried
neighbor cannot lie strictly inside: the edge to it would cross the
circuit without meeting a circuit vertex.  Planarity forbids this, and
expansion stops at vertices of the circuit because they have color
\(s\).  The positive separation of the narrower band from the protected
disk proves the asserted implication from a circuit to success.

For a topological rectangle, seed the exploration on the inspecting
side and expand color \(-s\) within the rectangle.  A color-\(-s\) path
reaching the opposite side certifies failure.  Otherwise the boundary
of the reached color-\(-s\) set facing that side gives a color-\(s\)
crossing between the two end sides.  To check this finite planar
assertion directly, join the reached components by a thin deterministic
strip immediately outside the inspecting side.  Trace the boundary of
a small polygonal neighborhood of this connected set on the side facing
the opposite boundary.  Since the reached set does not meet that
boundary, this trace has a component running between the two end sides.
Every graph vertex encountered just beyond the reached set was queried
and has color \(s\), because otherwise it would have been added to the
queue.  As the trace passes through a triangular face, consecutive such
vertices coincide or are joined by the remaining edge of that face.
They consequently give a walk of color \(s\) between the end sides.
Erasing closed subwalks gives a simple crossing.  Conversely a crossing
between the end sides blocks any opposite-color path between the
inspecting and opposite sides, since the two planar paths would have
to share a vertex.
Boundary arcs may be separated at their corners by one lattice step;
the prescribed buffered gates allow this harmless choice of finite
rectangle convention.  Thus the existence of the desired crossing is
determined by the transcript, and on success no queries pass strictly
beyond its facing boundary.  In applying this construction, the two end
gates lie past the comparison region, so a nearest-neighbor path cannot
go around a gate within that region.

Finally, a nearest-neighbor bond cannot cross a simple path or circuit
of the triangulation in the interior of an edge.  If its ends lie on
opposite sides, it must meet a path vertex.  This proves the incidence
assertion even when the separating path uses diagonal steps.
\end{proof}

Only the facing pins enter Proposition~\ref{prop:strict-incidence}.
Examined spins on the other side of a shield may have arbitrary signs,
and all of them remain part of the conditional law.  Likewise, when a
crossing test uses vertices outside one of its two comparison domains,
those vertices may be filled virtually by the extremal color for the
purpose of testing the increasing or decreasing event.  Such a filling
does not change any physical pin or any interaction.  The comparison
applies to the common free vertices after its incidence hypotheses have
been checked.

\subsection{From an extreme odds ratio to all interior densities}

Our final comparison will cut the auxiliary-bit graph along a square
inside the free buffer.  Conditional on all bit channels crossing that
cut, the interior law is the same in the domain and on the torus; the
exterior data change only the mixing law on the cut.  We must therefore
bound how much any two cut assignments can change the interior density.

For the bits of the \emph{unconditioned} torus, MTP\(_2\) reduces this
task to a single extreme odds ratio.  Two nested monochromatic shields
will bound that ratio using the incidence comparison and the circuit
probabilities from Section~\ref{sec:stability}.  We begin with the finite
algebraic reduction.  The four-difference identity below is the one used
in Lemma~2.1,
``Extreme rectangles in a ferromagnetic marginal'', in Section~2 of
\cite{OpenAIStatic2026}; its proof applies to any positive MTP\(_2\) mass.

\begin{lemma}[Rectangles of an MTP\(_2\) mass]
\label{lem:comparison-rectangles}
Let \(p(i,j)>0\) be an MTP\(_2\) probability mass on two disjoint finite
sets of binary coordinates \(I,J\).  Let \(i_-,i_+,j_-,j_+\) denote their
constant minus and plus configurations, set \(f=\log p\), and put
\begin{equation}
 \Lambda=f(i_+,j_+)+f(i_-,j_-)
              -f(i_+,j_-)-f(i_-,j_+)\ge0.
 \label{eq:comparison-extreme-ratio}
\end{equation}
Then, for arbitrary configurations \(a,a'\) on \(I\) and \(b,b'\) on
\(J\),
\begin{equation}
 \bigl|f(a,b)+f(a',b')-f(a,b')-f(a',b)\bigr|\le\Lambda.
 \label{eq:comparison-all-rectangles}
\end{equation}
Consequently every two conditional laws of \(I\) given \(J\), and every
two mixtures of those conditional laws, have pointwise density ratio
between \(e^{-\Lambda}\) and \(e^{\Lambda}\).
\end{lemma}

\begin{proof}
For \(u\le v\) and \(s\le t\), define
\[
 \Delta(u,v;s,t)=f(u,s)+f(v,t)-f(u,t)-f(v,s)\ge0.
\]
Writing \(D=f(a,b)+f(a',b')-f(a,b')-f(a',b)\), expansion gives the exact
identity
\begin{align*}
 \Lambda-D={}&\Delta(i_-,a';j_-,b)
             +\Delta(a,i_+;b',j_+)\\
             &+\Delta(a',i_+;j_-,b')
             +\Delta(i_-,a;b,j_+).
\end{align*}
Every term is nonnegative; no order between \(a,a'\), or between
\(b,b'\), is required.  Interchanging \(b,b'\) proves the other bound.

If \(q_b,q_{b'}\) are the two conditional masses on \(I\),
\eqref{eq:comparison-all-rectangles} says that
\(\log(q_b/q_{b'})\) has oscillation at most \(\Lambda\).  Its density
ratio has mean one under \(q_{b'}\), so its minimum is at most one and
its maximum is at least one.  Both bounds in the statement follow.
They persist after taking arbitrary mixtures in either law.
\end{proof}

We will also use preservation of MTP\(_2\) under marginalization.  Here
is a short proof in the finite setting at hand.  Sum one binary
coordinate \(z\) out of a positive MTP\(_2\) mass \(h\), and check the
increment ratio in another coordinate \(i\).  That ratio for the summed
mass is the conditional average, under \(h\) with \(i=-1\), of
\[
 r(z,y)=\frac{h(i=+1,z,y)}{h(i=-1,z,y)},
\]
where \(y\) denotes the other coordinates.  MTP\(_2\) makes \(r\)
nondecreasing in \(z\) and in every coordinate of \(y\).  The
conditional probability of \(z=+1\), with \(i=-1\), also increases when
one of those coordinates of \(y\) increases, by the corresponding
two-coordinate inequality for \(h\).  Thus the averaged ratio
increases.  This verifies each two-coordinate square for the marginal.
Iterating proves the assertion.  In particular, summing the physical
spins from the unconditioned representation gives the bit mass
\(\prod_vF_v\), which is MTP\(_2\), and any marginal of this bit mass
is MTP\(_2\) as well.

\begin{proposition}[Interior comparison across a fixed buffer]
\label{prop:interior-density}
Fix a square \(Q\) and a larger square \(\widetilde Q\), with
\(\overline Q\subset\operatorname{int}\widetilde Q\), in a fixed
physical torus chart.  Take lattice mesh \(\delta\downarrow0\), with
torus side \(N\) chosen so that this chart has a fixed macroscopic
buffer from its translates.  Suppose every lattice vertex in
\(\widetilde Q\) is free in a domain spin law
\(\nu_{\Omega,\lambda}^{\tau}\).  For all sufficiently small mesh there
is a constant \(C<\infty\), depending on the two fixed squares and the
chart but not on \(\Omega,\tau\), or the mesh, such that
\begin{equation}
 C^{-1}\le
 \frac{\nu_{\Omega,\lambda}^{\tau}(\sigma_{Q\cap\delta\Z^2}=\eta)}
      {\mu_{\lambda,N}(\sigma_{Q\cap\delta\Z^2}=\eta)}
 \le C
 \quad\text{for every spin configuration }\eta.
 \label{eq:comparison-interior-density}
\end{equation}
The same bounds hold for every event determined by these spins.
The exterior pins may be arbitrary, including those induced by any
admissible exterior contour configuration.
\end{proposition}

\begin{proof}
All sets in this proof are chosen in physical coordinates before the
mesh tends to zero.  We identify a lattice vertex with its rescaled
position when using \(Q\) and \(\widetilde Q\).  Let \(I\) contain all
bits incident to the vertices of \(Q\).  Choose two surrounding square
annular bands \(A_1,A_2\), with \(A_1\) inside \(A_2\), and then a
square cut outside \(A_2\), all within \(\widetilde Q\).  Leave positive
distances between successive objects.  The inner protected disk for an
exploration in \(A_1\) contains the endpoints of every bit in \(I\);
the inner protected disk for an exploration in \(A_2\) contains all of
\(A_1\).  Choose slightly wider exploration bands around \(A_1,A_2\),
still disjoint and separated from the cut.  A circuit in the original,
narrower band then implies exploration success by
Lemma~\ref{lem:comparison-stopped-reveal}.

More precisely, take a lattice square vertex set \(S\) containing both
exploration bands, with its boundary separated from them, and let \(J\)
be \emph{all} bit channels on the
nearest-neighbor bonds with one endpoint in \(S\) and the other outside.
We require \(I\) to lie strictly inside this cut.  Let
\(\mathcal B_{\mathrm{int}}(S)\) be the bits on bonds with both endpoints
in \(S\).  Every physical vertex in \(S\) is free.  Once \(J=j\) is
fixed, the conditional mass of the inside spins and bits is
\[
 K_S(\sigma_S,b_{\mathcal B_{\mathrm{int}}(S)}\mid j)
 =\frac1{Z_S(j)}\prod_{v\in S}W_v(\sigma_v,b_{\mathcal B(v)}),
 \qquad b_J=j.
\]
Every bit incident to a vertex of \(S\) is either an interior bit or
belongs to \(J\); all remaining vertex factors lie outside \(S\).
This proves the factorization and shows that \(K_S\) is identical in
the domain and torus laws.  In particular, their inside-bit marginals
can differ only through their mixing distribution on \(j\).
The representation is the same for all the laws by
Proposition~\ref{prop:bit-representation}; exterior pins in the spin
model are retained exactly by Lemma~\ref{lem:contour-spin}.

Let \(\widehat\mu\) denote the joint torus law of spins and bits, and
write \(\widehat\mu_j=\widehat\mu(\,\cdot\mid J=j)\).  For a sign
\(s\), let \(C_r^s\) be the fixed spin test of a color-\(s\) circuit in
\(A_r\), and let \(E_r^s\) be success of the exploration with that
color in its wider band.  These are different events.  The geometric
implication is
\begin{equation}
 C_r^s\subseteq E_r^s.
 \label{eq:comparison-circuit-success}
\end{equation}
Only the fixed tests \(C_r^s\) are used as inputs to bulk stability.
By Theorem~\ref{thm:bulk-stability} and the pure simultaneous
signed-circuit positivity of Proposition~\ref{prop:pure-rsw}, for all
sufficiently fine meshes
\begin{equation}
 \widehat\mu(C_1^+\cap C_2^-)
 \ge c,\qquad
 \widehat\mu(C_1^-\cap C_2^+)\ge c
 \label{eq:comparison-two-circuit-positivity}
\end{equation}
for a fixed \(c>0\).  We may take \(c\le1\).

We first show that, for \emph{every} cut configuration \(j\),
\begin{equation}
 \widehat\mu_j(C_1^+)\ge c,\qquad
 \widehat\mu_j(C_1^-)\ge c.
 \label{eq:comparison-circuits-all-cuts}
\end{equation}
Condition the unconditioned torus experiment on \(E_2^-\), and then on
a successful complete exploration transcript.  All physical spins of
\(A_1\) remain free.  Let \(V\) be the unqueried nearest-neighbor
component containing the protected inner disk, and hence all of \(A_1\).
It is a component of the common free set in the revealed experiment
and in \(\widehat\mu_j\).  Every neighbor outside \(V\) is a revealed
minus spin in the first experiment and a free spin in the second.
The queried outer boundary of the exploration separates \(V\) from
the cut, so no bit of \(J\) touches \(V\).  The minus-pin/free case of Proposition~\ref{prop:strict-incidence} orders
the conditional experiment below \(\widehat\mu_j\) on the spins of
\(A_1\).  Arbitrary further exterior bit conditions are allowed by that
proposition.  Average over the successful transcripts.  Since
\(C_1^+\) is increasing, this gives
\begin{align*}
 \widehat\mu_j(C_1^+)
 &\ge \widehat\mu(C_1^+\mid E_2^-)\\
 &=\frac{\widehat\mu(C_1^+\cap E_2^-)}
          {\widehat\mu(E_2^-)}
 \ge \widehat\mu(C_1^+\cap C_2^-)\ge c.
\end{align*}
The denominator is positive by
\eqref{eq:comparison-two-circuit-positivity}.  Interchanging the colors
and applying the comparison to the decreasing event \(C_1^-\) proves
the second inequality.  This step uses no monotonicity in the cut
configuration \(j\).

Now take the two extreme cut configurations \(j_-,j_+\), and denote the
corresponding laws by \(\widehat\mu_-\) and \(\widehat\mu_+\).
Under \(\widehat\mu_-\), the event \(E_1^+\) has probability at least
\(c\), by \eqref{eq:comparison-circuits-all-cuts} and
\eqref{eq:comparison-circuit-success}.  Conditional on any successful
transcript, the component containing all endpoints of \(I\) has facing
plus pins.  Compare it to \(\widehat\mu_+\) before any reveal.  The
plus-pin/free case of Proposition~\ref{prop:strict-incidence} makes
the former law larger on the bits \(I\).  The cut bits are outside the
comparison component in both experiments.  The increasing event
\(\{I=i_+\}\) therefore satisfies
\begin{equation}
 \widehat\mu_-(I=i_+)\ge
 \widehat\mu_-(E_1^+)\widehat\mu_+(I=i_+)
 \ge c\,\widehat\mu_+(I=i_+).
 \label{eq:comparison-plus-extreme}
\end{equation}
The opposite construction, starting from \(\widehat\mu_+\) and
revealing a minus shield, gives
\begin{equation}
 \widehat\mu_+(I=i_-)\ge
 c\,\widehat\mu_-(I=i_-).
 \label{eq:comparison-minus-extreme}
\end{equation}
These inequalities compare possibly very small probabilities; they do
not require a lower bound for either all-equal configuration on \(I\).

The joint torus bit marginal \(p(I,J)\) is positive and MTP\(_2\), by
Lemma~\ref{lem:comparison-margins} and marginalization as proved above.
Its extreme ratio satisfies
\[
 e^{\Lambda}
 =\frac{\widehat\mu_+(I=i_+)\widehat\mu_-(I=i_-)}
        {\widehat\mu_-(I=i_+)\widehat\mu_+(I=i_-)}
 \le c^{-2},
\]
by \eqref{eq:comparison-plus-extreme} and
\eqref{eq:comparison-minus-extreme}.  Apply
Lemma~\ref{lem:comparison-rectangles}.  It bounds the density ratio of
every two mixtures of the laws of \(I\) given \(J\) between \(c^2\)
and \(c^{-2}\).  One mixture is the torus law, and the other is the
domain law, by the exact cut factorization.

Finally, conditional on \(I\), the physical spins in \(Q\) have the
same product of Bernoulli kernels with plus probabilities \(k_v(I)\)
in both laws.  Integrating the pointwise bit-density inequalities
against this positive kernel gives
\eqref{eq:comparison-interior-density}, with \(C=c^{-2}\).
Summation gives the event version.  Every geometric separation was
fixed before taking the fine-mesh threshold, so the constant is
independent of the domain and its exterior pins.
\end{proof}

The buffer in Proposition~\ref{prop:interior-density} stays fixed when a
smaller test box shrinks inside \(Q\).  Consequently its constant can
be used for all boxes in a fine grid inside a compact chart.  Combining
it with the torus four-passage estimate of
Theorem~\ref{thm:bulk-stability} and
Proposition~\ref{prop:pure-four-passages} gives the compact-chart
bottleneck bound used in Section~\ref{sec:curves}.  That section gives
the grid argument with the order of the mesh and radius limits explicit.
